\documentclass[12pt, a4paper]{article}
\usepackage{graphicx}
\usepackage{geometry}
\usepackage[utf8]{inputenc}
\graphicspath{{./images}}

\geometry{top=3cm}

\title{KWAME NKRUMAH UNIVERSITY OF SCIENCE AND TECHNOLOGY\\
COLLEGE OF ENGINEERING\\
DEPARTMENT OF COMPUTER ENGINEERING\\ 
\includegraphics[scale=0.5]{{/knust.jpg}} 
\\ALLOCATION OF SENSORS TO MONITOR VITAL SOIL AND ATMOSPHERIC FACTORS THAT AFFECT PLANT GROWTH\\}


\begin{document}
\maketitle
\thispagestyle{empty}
\newpage
\pagenumbering{arabic}
\section*{Project Members}
\begin{itemize}
\item Name: Amoah Junior Albert\\ Index No.: 3022620\\ Student Id.: 20742304
\item Name: Bonney Samuel Nii Awuley\\ Index No.: 3028320\\ Student Id.: 20771143
\item Name: Lawson Isaac Nii Lante\\ Index No.: 3033820\\ Student Id.: 20746869 
\item Name: Akaya Uche Daniel\\ Index No.:8252419\\ Student Id.: 20451546
\end{itemize}

\newpage
\section*{\textbf{Abstract}}
Agricultural productivity and sustainability face significant challenges due to unpredictable environmental conditions and inefficient resource management. This project presents an innovative farm monitoring system utilizing Internet of Things (IoT) technology to address these issues. By deploying multiple sensor nodes across the farm, the system monitors critical environmental parameters: atmospheric temperature, humidity, soil pH, and soil moisture.

The sensor nodes form a mesh network, allowing direct communication between nodes and enhancing network reliability and scalability. The collected data is transmitted through this mesh network to an MQTT broker, which facilitates seamless data communication with cloud services on the Google Cloud Platform (GCP). There, advanced analytics and machine learning models generate actionable insights, empowering farmers with predictive analytics and real-time decision-making tools.

A user-friendly mobile app interface provides farmers with instant access to their farm's conditions, detailed analytics over various timeframes (daily, weekly, monthly, yearly), and critical alerts for immediate action.

The goals of this system are to optimize farming operations, increase crop yields, and promote sustainable agricultural practices. By providing farmers with precise and timely data, the system addresses the problem of inefficient resource use and environmental unpredictability. Preliminary results have shown a marked improvement in resource management and crop health, highlighting the system's potential to revolutionize modern farming.

This integrated solution represents a significant step forward in precision agricultural technology, offering a scalable and effective method to enhance productivity and sustainability in farming.

\end{document}